\documentclass{beamer}

\usetheme{Madrid}

%================= LANGUAGE =================
\usepackage[utf8]{inputenc}
\usepackage[T5]{fontenc}
\usepackage[vietnamese]{babel}

%================= PACKAGE =================
\usepackage{amsmath}
\usepackage{tikz}
\usepackage{listings}
\usepackage{xcolor}
\usepackage{tcolorbox}
\usepackage{fontawesome5}

%================= COLOR =================
\definecolor{main1}{RGB}{102,126,234}
\definecolor{main2}{RGB}{118,75,162}

\setbeamercolor{frametitle}{fg=white,bg=main1}
\setbeamercolor{title}{fg=white,bg=main2}

%================= BACKGROUND =================
\setbeamertemplate{background canvas}{
\begin{tikzpicture}[remember picture,overlay]
\path[left color=main1!20,right color=main2!20]
(current page.south west) rectangle (current page.north east);
\end{tikzpicture}
}

%================= BOX =================
\tcbset{
  colback=white,
  colframe=main1,
  arc=3mm,
  boxrule=0.8pt
}

%================= CODE =================
\lstset{
  language=C++,
  basicstyle=\ttfamily\tiny,
  keywordstyle=\color{main2},
  commentstyle=\color{green!50!black},
  breaklines=true
}

%================= INFO =================
\title{Thuật toán BFS (Breadth-First Search)}
\author{Lê Văn Hoàng An}
\date{\today}

\begin{document}

\frame{\titlepage}

%================= INTRO =================
\section{Giới thiệu}

\begin{frame}{Đồ thị là gì?}
\begin{tcolorbox}
\begin{itemize}
\item \faProjectDiagram Đồ thị gồm tập đỉnh (Vertex) và cạnh (Edge)
\item \faShareAlt Biểu diễn mối quan hệ giữa các đối tượng
\item \faRandom Có thể là đồ thị vô hướng hoặc có hướng
\item \faGlobe Ứng dụng: mạng xã hội, bản đồ, AI, hệ thống gợi ý
\end{itemize}
\end{tcolorbox}

\vspace{0.2cm}
\begin{tcolorbox}
\textbf{Ví dụ thực tế:}
\begin{itemize}
\item Facebook: người dùng = đỉnh, kết bạn = cạnh
\item Google Maps: địa điểm = đỉnh, đường đi = cạnh
\end{itemize}
\end{tcolorbox}
\end{frame}

%================= GRAPH =================
\section{Minh họa}

\begin{frame}{Minh họa đồ thị và BFS}
\begin{columns}

\column{0.5\textwidth}
\begin{tikzpicture}[scale=1.1, every node/.style={circle, draw}]
\node (A) at (0,0) {A};
\node (B) at (2,1) {B};
\node (C) at (2,-1) {C};
\node (D) at (4,1) {D};
\node (E) at (4,-1) {E};

\draw (A)--(B);
\draw (A)--(C);
\draw (B)--(D);
\draw (C)--(E);
\end{tikzpicture}

\column{0.5\textwidth}
\begin{tcolorbox}
\begin{itemize}
\item \faPlay Bắt đầu từ A
\item \faLayerGroup BFS duyệt theo từng lớp
\item \faArrowRight A → B, C
\item \faArrowRight Tiếp → D, E
\end{itemize}
\end{tcolorbox}

\end{columns}
\end{frame}

%================= THEORY =================
\section{Lý thuyết}

\begin{frame}{Ý tưởng + Giả mã BFS}
\begin{tcolorbox}
\begin{itemize}
\item \faBullseye Mục tiêu: duyệt theo từng lớp (level)
\item \faDatabase Cấu trúc: Queue (FIFO)
\item \faCheckCircle Mỗi đỉnh thăm đúng 1 lần
\end{itemize}
\end{tcolorbox}

\vspace{0.2cm}

\begin{tcolorbox}
\textbf{Quy trình thực hiện:}
\begin{enumerate}
\item Khởi tạo:
\begin{itemize}
\item Đánh dấu tất cả là chưa thăm
\item Đưa đỉnh start vào queue
\end{itemize}

\item Lặp:
\begin{itemize}
\item Lấy đỉnh đầu $u$
\item Xử lý $u$
\item Duyệt các đỉnh kề $v$
\item Nếu chưa thăm → thêm vào queue
\end{itemize}
\end{enumerate}
\end{tcolorbox}

\vspace{0.2cm}
\begin{tcolorbox}
\textbf{Độ phức tạp:} $O(V + E)$
\end{tcolorbox}
\end{frame}

%================= CODE =================
\section{Cài đặt}

\begin{frame}[fragile]{Code BFS (C++)}
\begin{tcolorbox}
\begin{lstlisting}
#include <bits/stdc++.h>
using namespace std;

void BFS(int start, vector<vector<int>>& adj) {
    vector<bool> visited(adj.size(), false);
    queue<int> q;

    visited[start] = true;
    q.push(start);

    while (!q.empty()) {
        int u = q.front();
        q.pop();

        // Xu ly dinh u
        cout << u << " ";

        for (int v : adj[u]) {
            if (!visited[v]) {
                visited[v] = true;
                q.push(v);
            }
        }
    }
}

/*
- Thoi gian: O(V + E)
- Duyet theo tung lop
- Dung trong tim duong ngan nhat
*/
\end{lstlisting}
\end{tcolorbox}
\end{frame}

%================= EXAMPLE =================
\section{Ví dụ}

\begin{frame}{Ví dụ minh họa BFS}
\begin{columns}

\column{0.5\textwidth}
\begin{tikzpicture}[scale=1.1, every node/.style={circle, draw}]
\node (1) at (0,0) {1};
\node (2) at (2,1) {2};
\node (3) at (2,-1) {3};
\node (4) at (4,1) {4};
\node (5) at (4,-1) {5};

\draw (1)--(2);
\draw (1)--(3);
\draw (2)--(4);
\draw (3)--(5);
\end{tikzpicture}

\column{0.5\textwidth}
\begin{tcolorbox}
\begin{itemize}
\item \faRoute Thứ tự:
\[
1 \rightarrow 2 \rightarrow 3 \rightarrow 4 \rightarrow 5
\]
\item \faLayerGroup Lớp:
\begin{itemize}
\item L0: 1
\item L1: 2, 3
\item L2: 4, 5
\end{itemize}
\end{itemize}
\end{tcolorbox}

\end{columns}
\end{frame}

%================= ANIMATION =================
\section{Animation BFS}

\begin{frame}{Minh họa BFS từng bước}
\centering

\begin{tikzpicture}[
    scale=1.2,
    every node/.style={circle, draw, minimum size=8mm},
    visited/.style={fill=green!40},
    current/.style={fill=yellow!60}
]

\node (1) at (0,0) {1};
\node (2) at (2,1) {2};
\node (3) at (2,-1) {3};
\node (4) at (4,1) {4};
\node (5) at (4,-1) {5};

\draw (1)--(2);
\draw (1)--(3);
\draw (2)--(4);
\draw (3)--(5);

\only<1>{\node[current] at (1) {1};}
\only<2>{\node[visited] at (1) {1}; \node[current] at (2) {2}; \node[current] at (3) {3};}
\only<3>{\node[visited] at (1) {1}; \node[visited] at (2) {2}; \node[visited] at (3) {3}; \node[current] at (4) {4};}
\only<4>{\node[visited] at (1) {1}; \node[visited] at (2) {2}; \node[visited] at (3) {3}; \node[visited] at (4) {4}; \node[current] at (5) {5};}
\only<5>{\node[visited] at (1) {1}; \node[visited] at (2) {2}; \node[visited] at (3) {3}; \node[visited] at (4) {4}; \node[visited] at (5) {5};}

\end{tikzpicture}

\vspace{0.5cm}
\only<1>{Bắt đầu từ đỉnh 1}
\only<2>{Thêm 2, 3 vào queue}
\only<3>{Xử lý 2 → thêm 4}
\only<4>{Xử lý 3 → thêm 5}
\only<5>{Hoàn tất BFS}

\end{frame}

%================= APPLICATION =================
\section{Ứng dụng}

\begin{frame}{Ứng dụng của BFS}
\begin{tcolorbox}
\begin{itemize}
\item \faRoute Tìm đường đi ngắn nhất (không trọng số)
\item \faProjectDiagram Kiểm tra liên thông
\item \faRobot AI, game, mê cung
\item \faUsers Mạng xã hội (gợi ý bạn bè)
\end{itemize}
\end{tcolorbox}
\end{frame}

%================= END =================
\begin{frame}{Kết luận}
\begin{tcolorbox}
\begin{itemize}
\item BFS là thuật toán nền tảng
\item Độ phức tạp tối ưu $O(V+E)$
\item Rất quan trọng trong thực tế
\end{itemize}
\end{tcolorbox}

\centering
\vspace{0.5cm}
\Huge{\textbf{Cảm ơn!}}
\end{frame}

\end{document}